\documentclass[11pt]{article}

\usepackage[margin=1in]{geometry}
\usepackage{amsmath,amssymb,amsthm,mathtools}
\usepackage[hidelinks]{hyperref}

\newtheorem{theorem}{Theorem}
\newtheorem{lemma}[theorem]{Lemma}
\newtheorem{corollary}[theorem]{Corollary}
\theoremstyle{remark}
\newtheorem{remark}[theorem]{Remark}

\newcommand{\cP}{\mathcal P}
\newcommand{\supp}{\operatorname{supp}}
\newcommand{\Sub}{\operatorname{Sub}}
\newcommand{\NU}{\operatorname{NU}}
\newcommand{\HJ}{\operatorname{HJ}}
\newcommand{\MHJ}{\operatorname{MHJ}}
\newcommand{\Var}{\operatorname{Var}}

\title{The Graham--Rothschild Theorem for Combinatorial Lines\\
from the Multidimensional Hales--Jewett Theorem}
\author{}
\date{}

\begin{document}
\maketitle

\section{Statement}

Let $A$ be a finite alphabet, $|A|=k\ge 2$, and let $v$ be a symbol not in
$A$. A \emph{variable word} of length $N$ is an element of
$(A\cup\{v\})^N$ in which $v$ occurs. It generates the combinatorial line
\[
  \{w(a):a\in A\}\subseteq A^N.
\]
More generally, an $m$-variable word is a word over
$A\cup\{v_1,\dots,v_m\}$ in which every $v_j$ occurs. We use the usual
block convention: all occurrences of $v_j$ precede all occurrences of
$v_{j+1}$. It generates an $m$-dimensional combinatorial subspace of $A^N$.

If $w=w(v_1,\dots,v_m)$ is an $m$-variable word, let $\Sub_1(w)$ be the set
of all variable words
\[
  w(\alpha_1,\dots,\alpha_m),
  \qquad
  \alpha_j\in A\cup\{v\},
\]
for which at least one $\alpha_j$ equals $v$. These are precisely the
combinatorial lines contained in the space generated by $w$.

\begin{theorem}[Graham--Rothschild for lines]\label{thm:GR-lines}
For every $k,m,r\ge 1$ with $k\ge 2$, there is an integer $N$ such that,
for every alphabet $A$ of cardinality $k$ and every $r$-coloring of the
variable words of length $N$ over $A$, there is an $m$-variable word $w$
of length $N$ for which $\Sub_1(w)$ is monochromatic.
\end{theorem}

We prove this from the ordinary and multidimensional colored
Hales--Jewett theorems, together with the elementary finite Ramsey theorem.
In fact, the multidimensional theorem with target dimension $1$ is the
ordinary theorem, so one may regard multidimensional Hales--Jewett as the
only Hales--Jewett input.

\section{Canonizing the constant letters}

The first lemma says that, after passing to a sufficiently large
combinatorial subspace, the color of a line depends only on which ambient
variables are active, and not on the constants assigned to the other
variables.

\begin{lemma}[support canonization]\label{lem:canonization}
For every $k,r,\ell\ge 1$ there is an integer $N$ with the following
property. If $A$ is an alphabet of cardinality $k$ and
\[
  c:\Var(A,N)\longrightarrow [r]
\]
is an $r$-coloring of the variable words of length $N$, then there is an
$\ell$-variable word
\[
  w=w(x_1,\dots,x_\ell)
\]
of length $N$ such that the color of
\[
  w(\alpha_1,\dots,\alpha_\ell),
  \qquad \alpha_i\in A\cup\{v\},
\]
depends only on the nonempty set
\[
  \supp(\alpha)=\{i\in[\ell]:\alpha_i=v\}.
\]
\end{lemma}

\begin{proof}
Put $A_v=A\cup\{v\}$. We first choose positive integers
$n_1,\dots,n_\ell$. Set $S_0=0$, and, after $n_1,\dots,n_{i-1}$ have been
chosen, let $S_{i-1}=n_1+\cdots+n_{i-1}$ and choose $n_i$ so large that
ordinary Hales--Jewett applies to the alphabet $A$ with
\[
  r^{(k+1)^{S_{i-1}+\ell-i}}
\]
colors. Put $S_i=S_{i-1}+n_i$ and $N=S_\ell$.

We now work backwards. Initially let $c_\ell=c$, viewed as a coloring of
the variable words of length $S_\ell$. Suppose that $1\le i\le\ell$ and
that $c_i$ is an $r$-coloring of the variable words of length
$S_i+\ell-i$. For each constant word $u\in A^{n_i}$, form its color profile
\[
  \Phi_i(u)
  =\bigl(c_i(puq)\bigr)_{(p,q)},
\]
where
\[
  p\in A_v^{S_{i-1}},\qquad q\in A_v^{\ell-i},
\]
and the pair $(p,q)$ ranges over those for which the concatenation $pq$
contains $v$. There are at most
$r^{(k+1)^{S_{i-1}+\ell-i}}$ possible profiles. By the choice of $n_i$,
Hales--Jewett gives a variable word $u_i(x_i)$ of length $n_i$ such that
\[
  \Phi_i\bigl(u_i(a)\bigr)
\]
is independent of $a\in A$.

Define $c_{i-1}$ on the variable words of length
$S_{i-1}+1+\ell-i$ by
\[
  c_{i-1}(pyq)=c_i\bigl(pu_i(y)q\bigr),
  \qquad y\in A_v.
\]
This is well-defined: if $y\in A$, then $pq$ already contains $v$, while
if $y=v$, then $u_i(v)$ contains $v$. Moreover, whenever $pq$ contains
$v$, the value
\[
  c_{i-1}(paq)
\]
is independent of $a\in A$.

After completing the backwards construction, set
\[
  w(x_1,\dots,x_\ell)
  =u_1(x_1)u_2(x_2)\cdots u_\ell(x_\ell).
\]
Suppose that $\alpha,\beta\in(A\cup\{v\})^\ell$ have the same nonempty
support. One may pass from $\alpha$ to $\beta$ by changing one constant
coordinate at a time. When the $i$th constant coordinate is changed, some
other coordinate still equals $v$, because the common support is nonempty.
The insensitivity obtained at stage $i$ therefore shows that the color does
not change. Hence the color depends only on the support.
\end{proof}

\section{Finite block unions from multidimensional Hales--Jewett}

For pairwise disjoint nonempty sets $E_1,\dots,E_t$, write
\[
  \NU(E_1,\dots,E_t)
  =\left\{\bigcup_{i\in I}E_i:
  \varnothing\ne I\subseteq[t]\right\}.
\]
We write $E<F$ when $\max E<\min F$.

We assume the following multidimensional colored Hales--Jewett theorem: for
every alphabet size $q$, dimension $d$, and number of colors $r$, there is
$\MHJ(q,d,r)$ such that every $r$-coloring of $[q]^M$, with
$M\ge\MHJ(q,d,r)$, has a monochromatic $d$-dimensional combinatorial
subspace.

\begin{lemma}[a translated finite-unions cube]\label{lem:translated-cube}
Let $n,r\ge1$ and put $M=\MHJ(2,n+1,r)$. For every $r$-coloring of the
nonempty subsets of $[M]$, there are pairwise disjoint nonempty sets
\[
  E,G_1,\dots,G_n\subseteq[M]
\]
such that all sets
\[
  E\cup\bigcup_{i\in I}G_i,
  \qquad I\subseteq[n],
\]
have the same color.
\end{lemma}

\begin{proof}
Extend the coloring arbitrarily to the empty set and identify
$\{0,1\}^M$ with $\cP([M])$ by characteristic functions. By
multidimensional Hales--Jewett, there is a monochromatic
$(n+1)$-dimensional subspace. Thus there are pairwise disjoint nonempty
wildcard sets
\[
  X_0,X_1,\dots,X_n
\]
and a fixed set $S$, disjoint from all the $X_i$, such that every set
\[
  S\cup\bigcup_{i\in J}X_i,
  \qquad J\subseteq\{0,1,\dots,n\},
\]
has the same color. Take
\[
  E=S\cup X_0,
  \qquad
  G_i=X_i\quad(1\le i\le n).
\]
All displayed sets in the statement correspond to points of the
monochromatic subspace with the zeroth wildcard set switched on.
\end{proof}

\begin{lemma}[minimum canonization]\label{lem:min-canonization}
For every $t,r\ge1$ there is an integer $D(t,r)$ such that every
$r$-coloring $d$ of the nonempty subsets of $[D(t,r)]$ admits pairwise
disjoint nonempty sets $E_1,\dots,E_t$ satisfying
\[
  d\left(\bigcup_{i\in I}E_i\right)
  \quad\text{depends only on }\min I
\]
for every nonempty $I\subseteq[t]$.
\end{lemma}

\begin{proof}
Induct on $t$. The case $t=1$ is immediate. Suppose the assertion is known
for $t$, and set
\[
  D(t+1,r)=\MHJ\bigl(2,D(t,r)+1,r\bigr).
\]
Apply Lemma~\ref{lem:translated-cube} with $n=D(t,r)$. We obtain a set $E$
and disjoint sets $G_1,\dots,G_{D(t,r)}$ such that
\[
  \{E\}\cup\{E\cup H:H\in\NU(G_1,\dots,G_{D(t,r)})\}
\]
is monochromatic. Color the nonempty subsets $I$ of $[D(t,r)]$ by
\[
  I\longmapsto d\left(\bigcup_{i\in I}G_i\right).
\]
By the inductive hypothesis, there are pairwise disjoint nonempty index
sets $I_2,\dots,I_{t+1}$ for which the color of their unions depends only on
the least selected index. Put
\[
  E_1=E,
  \qquad
  E_j=\bigcup_{i\in I_j}G_i\quad(2\le j\le t+1).
\]
A union containing $E_1$ has the color supplied by
Lemma~\ref{lem:translated-cube}; a union not containing $E_1$ has color
determined by the inductive hypothesis. This proves the claim.
\end{proof}

\begin{corollary}[finite disjoint unions]\label{cor:disjoint-unions}
For every $m,r\ge1$ there is an integer $D$ such that every $r$-coloring of
the nonempty subsets of $[D]$ admits pairwise disjoint nonempty sets
$E_1,\dots,E_m$ for which $\NU(E_1,\dots,E_m)$ is monochromatic.
\end{corollary}

\begin{proof}
Apply Lemma~\ref{lem:min-canonization} with
\[
  t=(m-1)r+1.
\]
Among the $t$ singleton colors $d(E_i)$, the pigeonhole principle gives
indices
\[
  i_1<\cdots<i_m
\]
with the same color. Every nonempty union of
$E_{i_1},\dots,E_{i_m}$ has the color of the member with least index, and
these $m$ colors are equal.
\end{proof}

For the block convention on variables, we need the following standard
block form. Its only additional input is the elementary finite Ramsey
theorem.

\begin{corollary}[finite block unions]\label{cor:block-unions}
For every $m,r\ge1$ there is an integer $L$ such that every $r$-coloring
\[
  d:\cP([L])\setminus\{\varnothing\}\longrightarrow[r]
\]
admits nonempty sets
\[
  \gamma_1<\gamma_2<\cdots<\gamma_m
\]
for which
\[
  d\left(\bigcup_{j\in J}\gamma_j\right)
\]
is independent of the nonempty set $J\subseteq[m]$.
\end{corollary}

\begin{proof}
Let $M$ be the number given by Corollary~\ref{cor:disjoint-unions} for
$(m,r)$. By finitely many applications of the finite Ramsey theorem, choose
$L$ so large that every $r$-coloring of the nonempty subsets of $[L]$ has
an $M$-element set
\[
  H=\{h_1<\cdots<h_M\}
\]
with the following property: for every $s\in[M]$, all $s$-element subsets
of $H$ have one common color, say $\rho_s$.

Color the nonempty subsets $I$ of $[M]$ by
\[
  I\longmapsto \rho_{|I|}.
\]
By Corollary~\ref{cor:disjoint-unions}, there are pairwise disjoint
nonempty sets $I_1,\dots,I_m\subseteq[M]$ such that
\[
  \rho_{|\bigcup_{j\in J}I_j|}
\]
is independent of nonempty $J\subseteq[m]$. Since the $I_j$ are disjoint,
this quantity is
\[
  \rho_{\sum_{j\in J}|I_j|}.
\]
Partition the first $|I_1|+\cdots+|I_m|$ points of $H$ into successive
blocks
\[
  \gamma_1<\cdots<\gamma_m,
  \qquad |\gamma_j|=|I_j|.
\]
Then
\[
  \left|\bigcup_{j\in J}\gamma_j\right|
  =\sum_{j\in J}|I_j|,
\]
so all these unions have the same original color.
\end{proof}

\section{Proof of the line case}

\begin{proof}[Proof of Theorem~\ref{thm:GR-lines}]
Fix $k,m,r$, and let $L$ be supplied by
Corollary~\ref{cor:block-unions}. Apply
Lemma~\ref{lem:canonization} with $\ell=L$, and let $N$ be the resulting
integer.

Let $c$ be an $r$-coloring of the variable words of length $N$. By
Lemma~\ref{lem:canonization}, there is an $L$-variable word
\[
  w=w(x_1,\dots,x_L)
\]
such that the color of a one-variable subword of $w$ depends only on its
support. Hence there is a well-defined coloring
\[
  \widehat c:\cP([L])\setminus\{\varnothing\}\longrightarrow[r]
\]
defined as follows: for nonempty $I\subseteq[L]$, let $\widehat c(I)$ be
the color of any word
\[
  w(\alpha_1,\dots,\alpha_L)
\]
with $\alpha_i=v$ exactly for $i\in I$ and $\alpha_i\in A$ otherwise.

Apply Corollary~\ref{cor:block-unions} to $\widehat c$. We obtain
\[
  \gamma_1<\gamma_2<\cdots<\gamma_m
\]
such that all nonempty unions of the $\gamma_j$ have one common
$\widehat c$-color. Fix a letter $a_0\in A$, and define an $m$-variable
subword $z=z(v_1,\dots,v_m)$ of $w$ by making the substitution
\[
  x_i\longmapsto
  \begin{cases}
    v_j,& i\in\gamma_j,\\
    a_0,& i\notin\gamma_1\cup\cdots\cup\gamma_m.
  \end{cases}
\]
Because $\gamma_1<\cdots<\gamma_m$, this is an $m$-variable word in block
position.

Consider any element of $\Sub_1(z)$. It has the form
\[
  z(\beta_1,\dots,\beta_m),
  \qquad \beta_j\in A\cup\{v\},
\]
where
\[
  J=\{j\in[m]:\beta_j=v\}
\]
is nonempty. Viewed as a one-variable subword of $w$, its support is
exactly
\[
  \bigcup_{j\in J}\gamma_j.
\]
The support-canonization lemma says that its $c$-color is
\[
  \widehat c\left(\bigcup_{j\in J}\gamma_j\right),
\]
and the block-unions corollary says that this value is independent of the
nonempty set $J$. Thus every member of $\Sub_1(z)$ has the same color.
\end{proof}

\begin{remark}
Via the unique correspondence between variable words and combinatorial
lines, Theorem~\ref{thm:GR-lines} says that every finite coloring of the
lines of a sufficiently high-dimensional word space $A^N$ has an
$m$-dimensional combinatorial subspace all of whose lines have the same
color. For a family $\mathcal L$ of lines, apply the theorem to the
$2$-coloring ``belongs to $\mathcal L$'' versus ``does not belong to
$\mathcal L$.''
\end{remark}

\begin{thebibliography}{9}

\bibitem{DK}
P.~Dodos and V.~Kanellopoulos,
\emph{Ramsey Theory for Product Spaces},
Chapter~2, especially the multidimensional Hales--Jewett theorem, the
disjoint-unions theorem, and the canonization of variable positions.

\bibitem{DKT}
P.~Dodos, V.~Kanellopoulos, and K.~Tyros,
\emph{A simple proof of the density Hales--Jewett theorem}.

\bibitem{Tyros}
K.~Tyros,
\emph{On colorings of variable words},
Discrete Math. 338 (2015), 1480--1483.

\end{thebibliography}

\end{document}
